% ---------------------------------------
%
%    Beispieldiplomarbeit
%
%    - text kann gel{\"o}scht, und mit eigenen Inhalten gef{\"u}llt werden. 
%	In dieser Datei kann die gesamte Konfiguration durchgeführt werden.*
%	
%
%    *mit Ausnahme der spezifischen Optionen in natbib.cfg, die nicht von Interesse sein sollten
% ---------------------------------------

\documentclass[
    smallheadings,  % kleinere {\"U}berschriften
    twoside,        % einseitig, nur rechte seiten
    liststotoc,     % listen in inhaltsverzeichnis aufnehmen
    bibtotoc,       % literaturverzeichnis in inhltsvz. aufnehmen
    headsepline,     % trennlinie unter kopfzeile
    12pt	%Schriftgröße
    ]{scrbook}

\usepackage{a4}  %a4 Seitenformat benutzen
\usepackage[left=4cm,right=2.5cm,top=3.7cm,bottom=4cm]{geometry}
\usepackage[english, ngerman]{babel} %Verwende deutsche und amerikanische Silbentrennung
\usepackage[utf8]{inputenc} %damit k{\"o}nnen Umlaute ganz normal geschrieben werden. 
\usepackage[numbers,square]{natbib}

\usepackage{listings} %Fuer Codelistings
\usepackage{subfigure} %f{\"u}r mehrteilige Grafiken
\usepackage{epsfig}    %damit funktioniert das einbinden von grafiken {\"u}ber epsfig.
\usepackage{graphicx}     % zum einbinden von grafiken
\graphicspath{{grafiken}{../}{kapitel}} %da sind m{\"o}gliche bilder fuer den includegraphics-Befehl zu finden (man muss dann nicht den ganzen Pfad bei includegraphics angeben. 

\usepackage{multirow}     %fuer kompliziertere Tabellen
\usepackage{longtable}	%fuer kompliziertere Tabellen
\usepackage{rotating}	%enable landscape pages
\usepackage{framed}
\usepackage{scrpage2}     % paket f{\"u}r kopf- und fu{\ss}zeilen
\pagestyle{scrheadings}   % kopzeilenseitenstil

%\usepackage{ifpdf} %provides the switch ``ifpdf'' to determine if latex or pdflatex is executed
%see ftp://ftp.ctan.org/tex-archive/macros/latex/contrib/oberdiek/ifpdf.pdf or doc folder

\usepackage{float} % adds location parameter [H] to [htb] to force figures and tables to a location, no floating

%%%% use hyperlinks and adjust settings, see http://en.wikibooks.org/wiki/LaTeX/Hyperlinks and
% ftp://tug.ctan.org/tex-archive/macros/latex/contrib/hyperref/doc/manual.pdf
\usepackage[
a4paper=true, % a4
pdftitle={Leveraging Continuous Integration process to improve Software Reuse}, %title to display for pdfs
pdfauthor={Piero A. Divasto Martínez}, %set author
%dvipdfmx, % must be used if pdf is created from the dvi file rather than by running pdflatex. Otherwise links won't work
colorlinks,% set link color for all links to black, i.e. invisible but working links
citecolor=black,%
filecolor=black,%
linkcolor=black,%
urlcolor=black, %
plainpages=false, %used with pdfpagelabels to make pdf readers show roman and arabic page numbers correctly
pdfpagelabels, %used with plainpages=false to make pdf readers show roman and arabic page numbers correctly
%pagebackref=true %false: no back-references in bibliography to citations in text
]{hyperref} %include hyperlinks for citations, urls, and references
%CAUTION: breaks color settings in dvi, pdf works fine
%%%%end: use hyperlinks

\usepackage{setspace}
\usepackage{url}          % fuer urls: schreibweise ist z.B. \url{http://www.uni-mannheim.de}
%\urlstyle{same} % makes urls in the same font as the rest of the text
\usepackage{xcolor}
\usepackage{minted}

\usepackage{titlesec}

%Automatisch Abkürzungsliste generieren: 2x latex, makeindex myfile.nlo -s nomencl.ist -o myfile.nls, 2x latex
\usepackage{nomencl} % um Abkürzungen aufzunehmen: \abbrev{PDA}{personal digital assistant}
\let\abbrev\nomenclature
\renewcommand{\nomname}{List of Abbreviations} %Titel der Liste hier ändern, eine von zwei Optionen nutzen
%\renewcommand{\nomname}{Abkürzungsverzeichnis} %Titel der Liste hier ändern, eine von zwei Optionen nutzen
\setlength{\nomlabelwidth}{.25\hsize}
\renewcommand{\nomlabel}[1]{#1 \dotfill}
\setlength{\nomitemsep}{-\parsep}
\makenomenclature
\newcommand{\Listofabbrev}{
\printnomenclature
\newpage
}

\usepackage{enumitem}


%Inhaltsverzeichnis
\usepackage
%[
%tocfullflat,			%alle Eintraege left-aligned, alternativ: tocflat
%tocbreakscareless		%page break between toc entry allowed
%]
{tocstyle}			%mehr Kontrolle ueber Inhaltsverzeichnis
\usetocstyle{allwithdot}	%ensure there are dots in table of contents, alternativ: noonewithdot, nopagecolumn
% docu http://www.tug.org/texlive/devsrc/Master/texmf-dist/source/latex/koma-script/tocstyle.dtx
\setcounter{secnumdepth}{3} %numbering up to fifth level in table of contents
\setcounter{tocdepth}{3} %show entries in toc up to fifth level (subsubsubsubsection)


\setlength{\parindent}{0pt}		% Einzug fuer neuen Absatz
\setlength{\parskip}\medskipamount % Abstand neuer Absatz, besser als explizite Angabe in pt

\onehalfspacing %Zeilenabstand 1,5


% kapitel{\"u}berschriften in schriftart mit serifen
\setkomafont{sectioning}{\normalfont\normalcolor\bfseries}

% gestaltung der kopfzeilen
\ohead{\pagemark}
\ifoot{}
\cfoot{}
\ofoot{} 
\cohead{}
\ihead{\headmark}
\setkomafont{pagehead}{\normalfont\bfseries}
\setkomafont{pagenumber}{\normalfont\bfseries}
\automark{section}

% ----- ende der pr{\"a}ambel ----------------------------------






\begin{document}  % dokument f{\"a}ngt an
\selectlanguage{english} %englische Silbentrennung, fuer deutsche Arbeiten: \selectlanguage{ngerman}
\frontmatter      % vorspann, kapitel r{\"o}misch nummeriert
\newgeometry{margin=3cm}
% Die Titelseite der Arbeit

\begin{titlepage}

\begin{center} % zentrieren

  % Logo der Universit{\"a}t Mannheim
  \begin{figure}[ht]
    \centering
    \includegraphics[width=.6\textwidth]{grafiken/unilogo}
  \end{figure}
  
  % Vertikaler Zwischenraum
  \bigskip
  \vfill 
  \begin{framed}
  % Titel der Arbeit und Typ der Arbeit, umrandet
    \begin{center}
     \textsc{{\Large Seizing Continuous Integration process to improve Software Reuse}}
                                % Letztes \\ ist wichtig, beginnt eine neue Zeile f{\"u}r die Art der Arbeit
  
      \bigskip
  
                                % Art der Arbeit, ggf. auszutauschen gegen Seminar- oder Doktorarbeit
      \textbf{Master Thesis}
    \end{center}
    \end{framed}
    \vfill
    \vfill
  
  % Daten des Erstellers, Einreichungsdatum
  % in einer Tabelle ausgerichtet
  \begin{tabular*}{0.62\textwidth}{r@{\extracolsep{\fill}}l}
   submitted: &\ February 2017\\\\
    by: &\ Piero Antonio Divasto Martínez\\
    &\ born March 10th 1986\\
    &\ in Iquique\\
    &\ Chile\\
    \\
    Student ID Number: &\ 1472793\\
  \end{tabular*}
  \vfill
  \vfill
  
  % Unten: Kontaktdaten des Lehrstuhls f{\"u}r Wirtschaftsinformatik 1
  
  \rule{\textwidth}{.4pt}\\ % vertikale Linie
  University of Mannheim\\
  Chair of Software Engineering\\
  B6, 26, Gebäudeteil C\\
  D -- 68159 Mannheim\\
  Phone: +49 621-181-1691, Fax +49 621-181-1692\\
  Internet: \url{http://swt.informatik.uni-mannheim.de/}
\end{center}

\end{titlepage} % Ende des Titelblatts

%%% Local Variables: 
%%% mode: latex
%%% TeX-master: "~/Documents/DA-Vorlage/beispiel/da-beispiel"
%%% End: 
     % titelseite einbinden
\restoregeometry
\chapter{Abstract}
On the one hand we have the Douglas McIlroy's vision \cite{McIlroy1968} of a software system composed of already existent components. Where you basically put different components, that already exist, together in order to form the system is being built. Although the IT industry has tried for many years to improve the speed and reduce cots of software development by reusing components, this vision is still the exception rather than the rule. On the other hand we have Continuous Integration which is a software development practice where members of a team integrate their work frequently, usually each person integrates at least daily, thus leading to multiple integrations per day. Each integration is verified by an automates build (including test) to detect integration errors as quickly as possible. In this work we present a novel approach that improve the reuse of software by leveraging continuous integration process. Also we present an unobtrusive tool as a proof-of-concept that extracts interface signatures and test classes from a given project, which are sent to a code-search engine to retrieve similar components based on text-driven and test-driven search techniques  % abstract einbinden
\tableofcontents            % inhaltsverzeichnis
\listoffigures              % abbildungsverzeichnis
\listoftables               % tabellenverzeichnis

%\clearpage %tell toc that there is a pagebreak after list of tables, however the hyperlink still points to the wrong page

\addcontentsline{toc}{chapter}{\nomname} %abkuerzungen ins inhaltsverzeichnis

% Abbreviations
\nomenclature{$CI$}{Continuous Integration}
\nomenclature{$VCS$}{Version Control System}
\nomenclature{$TDD$}{Test-driven Development}

\Listofabbrev % liste der abkuerzungen erstellen

\mainmatter       % hauptteil, kapitel lateinisch nummeriert
%\include{kapitel/content}
\selectlanguage{english} % jetzt sprechen wir wieder englisch
\chapter{How to use the template}
\label{chap:howto}

\textbf{Introduction}
\chapter{Software Reuse}
\label{chap:sw-reuse}

Software reuse is a concept that

\section{Citations}
\label{sec:citations}

The citation style used is called natbib and can be used to create numerical and author-year citations. To create the bibliography the bibliography style natdin is used as a basis for German papers, apalike2 for English ones. The respective .sty-file for natdin is located in the directory literatur. More information can be found on \url{http://merkel.zoneo.net/Latex/natbib.php} or in the directory ``doc''. Adjust the bib styple in the file arbeit.tex to fit your needs. Read the comments, some may help you. 

\subsection{The bib file}
\label{subsec:theBibFile}

The bibliography file is placed in the directory literatur and is called lit.bib. In the lit.bib file you will find different annotations to be used for different kinds of papers. This effects their representation in the bibliography. Some comments can be found in this file, too.

\subsection{Citations}
\label{subsec:citations}

Talking about a paper one may want to do a citation in the text using citet as in \citet{COMITY}. Alternatively parenthesis may be useful \citep{tannenbaum}. It is also possible to add more information like the pages of a book one refers to \citep[pp. 222-333]{tannenbaum}. \citet{ECORA}, \citet{random}, \citet{habil} and \citet{master} are in text citations. There is the Web Ontology Language \abbrev{OWL}{Web Ontology Language} \citep{OWL} as an example of an electronic resource.

There are three slightly different bib styles in the directory literatur. First, the native natdin.bst for DIN citations in German. Second, a slightly modified version ''natdinCustomized''. This style has minor changes in punctuation and is therefore not conform to DIN 1505. Third, an English translation of the latter is provided. Choose your bibliographystyle at the end of the root file (arbeit.tex). For English papers, apalike2 looks quite nice.

Several options for citations and the bibliography may be configured in the root file as usepackage parameters for natbib, e.g.\ the type of brackets used. See natnotes.pdf in the directory ``doc'' for more information. Additionally, for very detailed configuration it is possible to use the file ``natbib.cfg'' in the root directory to override predefined parameters.
\section{Continuous Integration}
\label{chap:ci}
In this chapter we explain about what Continuous Integration (CI) means, we describe the process steps, its benefits, challenges, and good practices of CI.

\subsection{What is Continuous Integration?}
\label{sec:ci-def}
The term can be found in the context of micro-processes development in the work of \citeauthor{Booch2007}. Then it was adopted by Kent Beck in his definition of Extreme Programming \cite{Beck1999}. However, it was Martin Fowler who was credited with establishing the current definitions of this practice. Fowler defines CI as following:

Continuous Integration is a software development practice where members of a team integrate their work frequently, usually each person integrates at least daily, thus leading to multiple integrations per day. Each integration is verified by an automates build (including test) to detect integration errors as quickly as possible. Many teams has found that this approach leads to significantly reduced integration problems and allows a team to develop cohesive software more rapid\cite{Fowler2006}.

CI emerges as a solution for the painful moment of software integration. Despite the fact that the process of integrating software is not a problem for a one-person project, when it increases in complexity it becomes more problematic. For instance, in the old days, a software was divided in modules and each of them were developed independently, once they done, those modules were put together in one step at the end of the project. Doing that led to all sorts of software quality problems, which are costly and often lead to project delays \cite{Duvall2007}.

\begin{figure}[ht]
	\centering
    \includegraphics[width=\textwidth]{grafiken/stress}
    \caption{Emotional cycle of manual delivery \cite{Goff2015}}
    \label{fig:emotional}
\end{figure}

Figure \ref{fig:emotional} describes the emotional cycle of a manual delivery process\cite{Goff2015} when the module integration is necessary. This step is tense moment as errors and failures appeared and they were difficult to find and fix at that stage of the development. In this manner, instead of waiting till the end of modules development to integrate, CI proposes to integrate frequently, usually one person should integrate at least once a day.

Therefore we can separte the CI workflow into the following steps \cite{Fowler2006}:

\textbf{CI Workflow}

\begin{itemize}
\item Developers check out code into their private workspaces
\item When done, commit the changes to the repository
\item The CI server monitors the repository and checks out changes when they occur.
\item The CI server builds the system and runs unit and integration tests.
\item The CI server releases deployable artifacts for testing.
\item The CI server assigns a build label to the version of the code it just built.
\item The CI server informs the team of the outcome of the build.
\item In case the build or test failed, the team fixes the issue at the earliest opportunity.
\item Continue to frequently integrate and test throughout the project.
\end{itemize}

This uncomplicated workflow helps the development team to focus on the current code change till it is verified and validated. Moreover, it provides continuous feedback on the quality of the code.

An extension of Continuous Integration in Continuous Delivery (CD). CD is a software development discipline where you build software in such a way that the software can be released to production at any time \cite{Fowler2013}. Figure \ref{fig:cicdpipeline} describes an example of a typical CI/CD pipeline where one can see the relationship between this two disciplines.

\begin{figure}[ht]
	\centering
    \includegraphics[width=\textwidth]{grafiken/cicd-pipeline}
    \caption{Typical CI/CD pipeline \cite{Chacon2009}}
    \label{fig:cicdpipeline}
\end{figure}

\subsubsection{CI Tools}

Although the CI is tool-agnostic, the selection of them depends on the project, framework in use, skill-set of the stakeholders and other factors. There are however two must-have tools of any CI system: (1) the Version Control System (VCS) and (2) the CI Server.

The most popular VCSs are SVN, Mercurial, and Git. On top of them, we can find Version Control Platforms (VCP) such as Bitbucket, Gitlab, and Github. In terms of CI servers, we can find Jenkins, Hudson, GoCD as open source projects; TravisCI, CircleCI, CodeShip and Team City as commercial tools.

Therefore choosing the appropriate tools is about finding the balance between price, setup, configuration efforts, ease-of-use, integration capabilities between the selected tools, framework suitability and maintainability in respect to current code base.

\textbf{Version Control Systems}

\paragraph{SVN}
Apache Subversion \footnote{https://subversion.apache.org/ (accessed: 13.01.2017)} is an open-source software versioning and revision control system, which is used to maintain current and historical versions of files such as source code, web pages, documents, etc. It appeared as an alternative to the Concurrent Versions System (CVS) \cite{Pilato2008}.

One can visualize the SVN filesystem as \emph{two-dimensional}. Two coordinates are used to unambiguously address filesystem items: (1) path, and (2) revision.

Each revision in a Subversion filesystem has its own root, which is used to access contents at that revision. Files are stored as links to the most recent change; thus a repository is rather compact. The system consumes storage space proportional to the number of changes made, not to the number of revisions. Moreover, the Subversion filesystem uses transactions to keep changes atomic. A transaction operates on a specified revision of the filesystem, not necessarily the latest. The transaction has its own root, on which changes are made. It is then either committed and becomes the latest revision, or is aborted. The transaction is actually a long-lived filesystem object; a client does not need to commit or abort a transaction itself, rather it can also begin a transaction, exit, and then can re-open the transaction and continue using it. Potentially, multiple clients can access the same transaction and work together on an atomic change, though no existing clients expose this capability \cite{Pilato2008}.
\paragraph{Mercurial}
Mercurial\footnote{https://www.mercurial-scm.org/ (accessed: 13.01.2017)} is a free, distributed source control management tool. It offers you the power to efficiently handle projects of any size while using an intuitive interface. It is easy to use and hard to break, making it ideal for anyone working with versioned files.

Mercurial's major design goals include high performance and scalability, decentralized, fully distributed collaborative development, robust handling of both plain text and binary files, and advanced branching and merging capabilities, while remaining conceptually simple.[3] It includes an integrated web-interface. Mercurial has also taken steps to ease the transition for users of other version control systems, particularly Subversion \cite{Mackall2006}.
\paragraph{Git}
Within the several different VCSs available, Git is one of the most popular\footnote{https://rhodecode.com/insights/version-control-systems-2016 (accessed: 14.01.2017)}. Git emerges in 2005 as an alternative to BitKeeper\footnote{https://www.bitkeeper.com/ (accessed: 13.01.2017)} after the break of the relationship between the commercial company behind BitKeeper and  the community that developed the Linux kernel \cite{Chacon2009}.

The main difference between Git and the other VCSs is the way it thinks about its data. Other VCSs such as Subversion, CVS, Perforce, and so on, think of the information they keep as a set of files and the changes made to each file over time. On the other hand, Git thinks its data like a set of snapshots of a miniature file system. Therefore, every time a user commits, or save the state of the project, Git basically takes a picture of what all the files look like at that moment and store a reference to that snapshot. In addition, Git has three main states where the files can reside in: committed, modified, and staged. Committed means that the data is stored in the local database. Modified means that the file was modified but it has not been committed yet. And staged means that the modified file has been marked to go into the next commit \cite{Chacon2009}. The figure \ref{fig:git-sates} depicts these three states.

\begin{figure}[ht]
	\centering
    \includegraphics[width=\textwidth]{grafiken/git-states}
    \caption{The three main states of Git \cite{Chacon2009}}
    \label{fig:git-sates}
\end{figure}

The git directory is the core of Git as it is where the metadata and object database for a project is stored. The working directory is a single checkout of one version of the project. Finally, the staging area is a file which is generally stored in the Git directory. This file contains information about what will go into the next commit.

\textbf{Version Control Platforms}

\paragraph{Bitbucket}
is a web-based hosting service for projects. It supports Mercurial and Git SCMs. Bitbucket\footnote{https://bitbucket.org (accessed: 21.02.2017)} offers both commercial plans and free accounts. It offers free accounts with an unlimited number of private repositories. Bitbucket is written in Python using the Django web framework. It started as an start-up by Jesper Nohr, but later it was acquired by Atlassian\footnote{https://www.atlassian.com/ (accessed: 21.02.2017)}.
\paragraph{Gitlab}
is an open-source web-based Git repository developed by GitLab Inc. The software was written mainly in Ruby by Dmitriy Zaporozhets and Valery Sizov. Although it offers, in general, the same features like Mercurial and Github, the main characteristics is that it allows you to deploy your own instance of GitLab in your own server.
\paragraph{Github}
is a web-based Git repository hosting service, which offers all the functionality of Git and its own features. It provides several collaboration tools such as wikis, bug tracking, feature requests, and task management as well as access control. Its development started in 2007 as a side project of P. J. Hyett and Chris Wanstrath \cite{Weis2014} and it was officially lunch on the 10th of April in 2008. Since then, Github\footnote{https://github.com/ (accessed: 13.01.2017)} has gain popularity through the community reaching more than 50 million projects hosted nowadays.\\

\textbf{Continuous Integration Servers}

\paragraph{Teamcity}
is a Java-based build management and continuous integration server from JetBrains\footnote{https://www.jetbrains.com (accessed: 21.02.2017)}. TeamCity\footnote{https://www.jetbrains.com/teamcity/ (accessed: 21.02.2017)} is commercial software and licensed under a proprietary license. It offers several features such as VCS interoperability, cloud integrations, code quality tracking, continuous integration, user management, system maintenance, build history, and extensibility and customization. 
\paragraph{TravisCI}
is a hosted, distributed continuous integration and continuous delivery service used to build and test software projects hosted in GitHub. It offers a free open-source\footnote{https://travis-ci.org/ (accessed: 21.02.2017)} and an enterprise version\footnote{https://enterprise.travis-ci.com/ (accessed: 21.02.2017)}. Currently more than 300.000 projects use TravisCI (free version).
\paragraph{Jenkins}
Jenkins is an open source automation server developed in Java. It was originally founded in 2006 as \textit{Hudson}, however a disputed with Oracle the community behind Hudson decided to change the project name to Jenkins\footnote{http://archive.is/fl45 (accessed: 13.01.2017)}.

Jenkins enables developers to reliable build, test, and deploy their software. Its extensible and plugin-based architecture has permitted to create a tons of plugins to adapt the CI server to a multitude of build, test, and deployment automation workloads. In 2015, Jenkins surpassed 100.000 known installation making it the most widely deployed automation server\footnote{https://jenkins.io/press/ (accessed: 13.01.2017)}.

In Table \ref{ci-comparison} we can see a comparison of the three CI Servers described above.

% Please add the following required packages to your document preamble:
% \usepackage{multirow}
% \usepackage{graphicx}
\begin{table}[]
\centering
\resizebox{\textwidth}{!}{%
\begin{tabular}{|c|c|l|l|l|l|l|l|c|c|c|}
\hline
\multirow{2}{*}{CI Server}     & \multirow{2}{*}{Platform}          & \multicolumn{3}{c|}{\multirow{2}{*}{License}}                                           & \multicolumn{1}{c|}{\multirow{2}{*}{Builders}}                                                                             & \multicolumn{1}{c|}{\multirow{2}{*}{Integrations}}                                                                                               & \multicolumn{1}{c|}{\multirow{2}{*}{IDE Support}}                                                                             & \multicolumn{3}{c|}{SCM Support}     \\ \cline{9-11} 
                               &                                    & \multicolumn{3}{c|}{}                                                                   & \multicolumn{1}{c|}{}                                                                                                      & \multicolumn{1}{c|}{}                                                                                                                            & \multicolumn{1}{c|}{}                                                                                                         & SVN        & Mercurial  & Git        \\ \hline
Teamcity                       & Web Container                      & \multicolumn{3}{l|}{Propietary}                                                         & \begin{tabular}[c]{@{}l@{}}MSBuild, NAnt, \\Visual Studio,  \\Ant, Maven 2-3, \\Gradle, Rake, \\command-line\end{tabular} & \begin{tabular}[c]{@{}l@{}}Jira, Bugzilla, \\FindBugs,  \\PMD, dotCover, \\NCover\end{tabular}                                                & \begin{tabular}[c]{@{}l@{}}Eclipse, Visual Studio, \\ IntelliJ IDEA, WebStorm, \\ PhpStorm, RubyMine, \\ PyCharm\end{tabular} & \checkmark & \checkmark & \checkmark \\ \hline
\multicolumn{1}{|l|}{TravisCI} & \multicolumn{1}{l|}{Hosted}        & \multicolumn{3}{l|}{MIT}                                                                & \begin{tabular}[c]{@{}l@{}}Ant, Maven, \\ Gradle\end{tabular}                                                              & Github, Heroku                                                                                                                                   & None                                                                                                                          & \texttimes & \texttimes & \checkmark \\ \hline
\multicolumn{1}{|l|}{Jenkins}  & \multicolumn{1}{l|}{Web Container} & \multicolumn{3}{l|}{\begin{tabular}[c]{@{}l@{}}CC\\ and MIT\end{tabular}} & \begin{tabular}[c]{@{}l@{}}MSBuild, NAnt, \\Ant, Maven, \\Kundo, Gradle,\\shell scripts,  \\ command-line\end{tabular}   & \begin{tabular}[c]{@{}l@{}}Bugzilla, Google Code, \\ Jira, Bitbucket, \\ Redmine,  \\ FindBugs, \\ Checkstyle, PMD, \\ Mantis, Trac\end{tabular} & \begin{tabular}[c]{@{}l@{}}Eclipse, IntelliJ IDEA, \\ NetBeans\end{tabular}                                                   & \checkmark & \checkmark & \checkmark \\ \hline
\end{tabular}%
}
\caption{Comparison of Continuous Integration servers}
\label{ci-comparison}
\end{table}

\subsection{Benefits in adopting CI/CD}
\label{ci-benefits}
Several are the benefits that come with the adoption of Continuous Integration in a software project. In this section we will describe five areas that are improved after CI implementation \cite{Rejstrom2016}. It is important to point out that these benefits come along with other practices such as agile transformation \cite{Laanti2011} and lean software development \cite{Poppendieck2003}.

\begin{itemize}
\item Shorter time-to-market: Many benefits come with the adoption of frequent releases. With fast and frequent releases is possible to get feedback quickly from customer and market, thus the organization gets better understanding of their needs expectations \cite{Neely2013}. With that is possible to focus on the most relevant features. Another benefit of delivering frequently is waste reduction as features can be deployed as soon as they are done \cite{Leppanen2015}. Furthermore, with frequent releases is possible to experiment with new features easily and with low impact \cite{Neely2013}.
\item Rapid feedback: With frequent releases is possible to show progress to customers and by that get feedback quickly. With that the development team can focus on important features instead of waste time in features that are not relevant for the customers or the market.
\item Improved software quality: Researches have reported a decreased in the number of open bugs and production incidents after adopting CI practices \cite{Mantyla2015} and the link between software quality improvement and the heavy reliance on automated test combined with smaller more manageable releases \cite{Leppanen2015}.
\item Improved release reliability: In the work of \cite{Neely2013} was proved that a working deployment pipeline along with intensive automated testing and fast rollback mechanism possitively affects release reliability and quality. With small and frequent releases fewer things can go wrong in a release according to \cite{Fowler2013}. In fact, reduction in the stress of developers and other stakeholder have been found by adopting CI \cite{Neely2013} \cite{Chen2015}.
\item Improved developer productivity: By automating deployment process, environment configuration and other non-value adding tasks, significant time saving for developers have been observed \cite{Rodriguez2016}. Also, in the work of \cite{Humble2010} was observed that although the setup cost of a deployment pipeline can be high, after it is setup developers can focus on value adding software development works.
\end{itemize}

\subsection{Challenges in adopting CI/CD}
\label{ci-challenges}

Adopting CI can be very beneficial for a software project as described above, however we can face some challenges that can jeopardize the successful implementation of it. According to the literature, there are seven common challenges in the implementation of a CI:

\begin{itemize}
\item Change resistance: Transforming the development of a project towards continuous integration practices requires investment and involvement from the entire organization \cite{Rodriguez2016}. Therefore any transformation in how an organization works, will receive a resistance on both a personal level and decision level within the organization.
\item External constraint: As a software project is part of a context, external constraint may appear. Normally customer preferences and domain imposed restrictions are sources of constraints. For example in highly restricted domains, legal regulations may require extensive testing before new version can be allowed to enter production \cite{Rodriguez2016}.
\item QA effort: The automated tests suit needs to be exhaustive enough in order to ensure the quality of what is being built. Thus it can lead to increase QA efforts due to difficulties in managing the test automation infrastructure \cite{Rodriguez2016}.
\item Legacy code: Normally legacy software has not been design for being automatically tested and may cause integration failures which may inhibit the continuous deployment process. The ownership of legacy code might belong to another company or team which shift the testing responsibility and this might delay the deployment process.
\item Complex software: If the complex if the software project is high, then setting up the CI workflow is more challenging \cite{Leppanen2015}.
\item Environment management: Keeping all the environments used in development, testing and production sync and similar can be challenging. This is due to the fact that differences in the environment configuration can lead to undetected issues appear in production. Therefore is essential to have a good configuration management that provisions environments automatically \cite{Leppanen2015}.
\item Manual testing: Although automated tests are very beneficial, some aspects of software need to be manually tested such as security issues, performance and UX/UI. Therefore 	heavy manual testing can impact the overall speed and smoothness of the process \cite{Leppanen2015}.
\end{itemize}

\subsection{Successful CI practices}
\label{ci-successful-pract}

In order to reduce the risks in adopting CI practices, eight principles, based on several literature reviews, were defined in \cite{Rejstrom2016} which should guide the CI implementation within organizations. These principles are generic solutions which can be adapted in most cases.

\begin{itemize}
\item Automate the build: The task that triggers the whole pipeline is the commit. After each commit, the next step should be build the binaries. The executable that result of the build, should go through the pipeline until it gets validated, verified and deployment ready.
\item Make your build self-testing: 
\item Every commit should build on an integration machine
\item Keep the build fast
\item Keep the build green
\item Test in a clone of the production environment
\item Everyone can see what is happening
\item Automate deployment
\end{itemize}

These 8 principles should guide the CI/CD implementation within organizations. The principles have been stablished as best practices through the validation of these concepts through several literature reviews [\cite{Rodriguez2016}; \cite{Mantyla2015}; \cite{Stahl2014}]. They can be viewed almost as a design pattern for adopting CI/CD, offering a generic solution adaptable to most cases.
\section{Microservices}

Microservices is a software architecture is usually linked to Martin Fowler in its article \cite{Fowler2014}. It is an architectural styles that aim to develop applications as a suite of small services independently of each other. Each of this services runs in its process and they communicate with lightweight mechanisms (e.g. HTTP resource API). These services are built around business capabilities and independently deployed by fully automated deployment machinery. Also there is a bare minimum of centralized management of these services, which may be written in different programming languages and use different data storage technologies \cite{Fowler2014}.

This architectural style emerges as a solution to the monolithic style due to its problems. An application built using monolithic architectural style is considered as a single unit with normally three main parts (figure \ref{}): a client-side user interface (i.e. HTML, CSS, JavaScript running on user's machine), a database and a server-side application. The latter is a monolith as it handles HTTP request, execute domain logic, retrieve and update data from the database, and populates and/or select the HTML views sent to the client-side. As it is a single logical executable, any changes made in the system, no matter the size of it, require a build and a deploy of the whole server-side application. It is hard to keep a good modular structure, making it harder to keep changes that must affect one module within that module. Also it is hard to scale as it is necessary to scale the entire application rather than parts of it requiring greater resources. These problems are causing frustrations on people that are using this style, specially nowadays when applications are deployed to the cloud, leading to the microservices architectural style.

\subsection{Characteristics of a Microservice Architecture}
In the work of Fowler and Lewis \cite{Fowler2014} described a list of common characteristic of applications built using microservice architectural style. It is important to point out that all the applications using this style must confirm every characteristic of this list.

\begin{description}[style=nextline]
\item[Componentization via Services] \hfill \
First of all it is important to explain the differences between libraries and services as the primary way of componentizing is by breking down software into services. On the one hand libraries are components that are linked into a application and called using in-memory function calls. On the other hand services are out-of-process components which communicates with a mechanism such as a web service request, or remote procedure call \cite{Fowler2014}.

Splitting a software into services brings benefits such as services as components are independently deployable making it easy for changes, or services have more explicit component interfaces by using explicit remote call mechanisms. 
\item[Organized around Business Capabilities] \hfill \
Normally in monolithic applications the development teams are divided by technology layers such as front-end team, server-side team and database team. However when team are divided in this way, even simple changes can lead to a cross-team project taking time and budgetary approval.

In Microservices approach this is taken in a different way by dividing the application into services organized around business capability. Thus services involve all needed for that business area, including user-interface (UI), database, business logic, and any other external collaboration. Therefore each teams are cross-functional, including the full range of skills required for the development of it.
\item[Products not Projects] \hfill \
Usually software development follows a project model where the objective is to deliver some piece of software which, when it is done, it is delivered to a maintenance team and the development team is disbanded.

Microservices approach follows the product model instead of project model. In the product model the development team takes full responsibility of the software, in other words it owns the product. This results in day-to-day contact of the developers with the users as they have to take on at least some of the support. Moreover this mentality ties in with the linkage to business capabilities as the focus is on how the software can assist its users to enhance the capabilities instead on put the focus on a list of functionality.
\item[Smart endpoints and dumb pipes] \hfill \
Microservices use the approach of smart endpoint and dumb pipes. Instead of putting so much smart into the communication mechanism itself, microservices approach aim to be as decoupled and as cohesive as possible. In other word, a microservice works as a black-box which receives a request, applies a logic, and produces a response. Usually two protocols are used, HTTP request-reponse with resource API's and lightweight messaging (e.g. RabbitMQ, ZeroMQ).
\item[Decentralized Governance] \hfill \
Using a centralized approach brings consequences such as standardise on single technology which keeps the team to use the right tool for the job. In contrast to that microservices approach allows to use use the technology that fit the requirements best. For instance, if NodeJS is more appropriate for a service than Java there will not problem choosing it.
\item[Decentralized Data Management] \hfill \
Microservices approach prefer letting each service manage its own database, either different instances of the same database technology, or entirely different database systems - an approach called Polyglot Persistence \cite{Leberknight2015}. 

Although this approach has some implications for managing updates or increment of complexity, the benefits are worth it. The common approach to solve the former implication is using transactions but implement it in services is difficult, that is why microservice architectures emphasize transactionless coordination between services, assuming that consistency may be only eventual consistency and problems are dealt with by compensation operations. This is useful in order to respond quickly to demand. In spite of the fact that by using different technologies increases the complexity, it allows the teams to solve the problem with the right tool. For instance if a service needs to just retrieve pages by looking up its ID, a key-value database is more suitable than a relational database. Moreover, as many NoSQL databases are designed to operate over clusters, they can tackle larger volumes of traffic and data which is harder with vertical scaling \cite{Fowler2011}.
\item[Design for failure] \hfill \
As a consequence of using microservice architecture, is that applications need to be designed so that they can tolerate the failure of services. As any service call could fail due to unavailability of the supplier service, the client service must respond as neat as possible.

Although this can be considered as a disadvantage in comparison to monolithic approach, using good monitoring and logging mechanism would help to overcome this. For instance a dashboard which display the status of the different services, current throughput, and latency would help to react in case one or more services go offline.
\end{description}

In conclusion, although microservices architectural style comes with many benefits, it also comes with several difficulties. By splitting up an application into services, we can simplify and speed up the release process, use the right tool for a job, or focus in how to assist better the user instead of a set of functionality. However we have to worry about changes to one service breaking its costumer, or deal with the increment of complexity.

It is important to consider that this style is not a silver bullet. As Fowler described in his blog post \cite{Fowler2014}, the core of The Guardian website \footnote{The Guardian: \url{https://www.theguardian.com/} (accessed: 04.02.2017).} was built as a monolith, however every new feature has been added as a microservice that uses the core's API.
\chapter{Usage Scenario}
\label{usage-scenario}
Before getting into the approach presented in this work, we will propose a usage scenario to strength and motivate it. The following scenario is based on the usage scenario found in \cite{Kessel2016}.

First of all we will assume that the test software project has the following characteristics. The project has adopted Continuous Integration as we described in Section \ref{chap:ci}, using Git as Version Control System, Jenkins as CI Server, and Java as programming language. Moreover the software is developed following the Test-driven development process.

Envisage the following scenario in which a development team is working or has code ownership of a web service application, and recognizes that she/he requires the ability to encode and decode information in the common Base64 format (e.g., the Base64 variant used in MIME transfer encoding \cite{rfc2045}).

As the project is being developed using TDD, developers create a test case that defines a function before it is implemented. For instance, if we follow the scenario, a member of the team will implement the method \emph{encode} which accept an array of bytes decoded in Base64 and it will return an array of chars which represents the input decoded. Moreover the method \emph{decode} that receives an array of chars and returns an array of bytes that represents the input encoded will need to be implemented too. However, as they are using TDD, this member must create a test of these methods before writing functional code. Listing \ref{Base64.java} represents the methods \emph{encode} and \emph{decode} and Listing \ref{Base64Test.java} represents the test class for these methods.

\lstinputlisting[
  language=Java, numbers=left, stepnumber=5, firstnumber=1, breaklines=true, 
  basicstyle=\footnotesize,
  numberstyle=\tiny,
  caption={Base64.java},
  captionpos=b,
  label=Base64.java
]
{code/Base64.java.txt}

\lstinputlisting[
  language=Java, numbers=left, stepnumber=5, firstnumber=1, breaklines=true, 
  basicstyle=\footnotesize,
  numberstyle=\tiny,
  caption={Base64Test.java},
  captionpos=b,
  label=Base64Test.java
]
{code/Base64Test.java.txt}

After that, this programmer will write enough code to fulfill the test case. This enough code neither need to be the \emph{best} nor the last version, because the next step will be to refactor it. Listing \ref{Base64.B.java} represents the enough code to pass the test.

\lstinputlisting[
  language=Java, numbers=left, stepnumber=5, firstnumber=1, breaklines=true, 
  basicstyle=\footnotesize,
  numberstyle=\tiny,
  caption={Base64.java with implementation to pass the test},
  captionpos=b,
  label=Base64.B.java
]
{code/Base64.java.B.txt}

Additionally, as the team adopted CI, developers should commit often the changes to the Git server. After adding enough functional code to fulfill the test case, the developer is able to push new code to the Git server. This action triggers the CI server which checkouts the code and runs the unit tests available. Therefore going along with the example given above, the class \emph{Base64.java} and the test class \emph{Base64Test.java} should be committed to the repository.

From this simple scenario emerge the following questions:

\begin{description}[style=nextline]
\item[How might the team realize that what they are working on has not been developed already?] \hfill \
If we follow the example, the implementation of encode and decode in Base64 is very easy to find, however that does not mean that the developer is aware of that at the moment he/she is developing. Therefore it would be very useful that an automated tool analyzes the code and recommends similar components that implement the functionality being developed.
\item[How might we seize continuous integration in order to improve chances of reusing software?] \hfill \
Continuous integration process seems to be a very good opportunity for software reuse. According to CI workflow explained in Section \ref{chap:ci}, page \pageref{sec:ci-def}, every change made in the code should be committed and pushed to the VCS. This pushed code is tested by the CI server and accepted if it passes all the tests otherwise it is rejected. This is a good chance to extract method signatures and make the code search named previously.
\item[Would TDD be helpful?] \hfill \
According to TDD process before implementing a method a failure test should be implemented before. If we assume that and that every code change is committed and pushed to the VCS, we will also be able to extract the test classes to search for similar components.
\end{description}

In the following Chapter we will answer these questions and present our approach that improve that chances of reusing software by seizing the use of Continuous Integration.
\chapter{Simile }
\label{chap:simile}

In this work we present Simile, which is a tool that leverages the Continuous Integration process in order to recommend similar components that are being developed in a software project. Following the usage scenario explained previously in section \ref{usage-scenario}, we will describe how our approach works:

Once a member of the team makes a commit and pushes to the remote Git server, our tool will be triggered by Jenkins. This tool will receives the repository URL where the project is being stored. Then it will clone the repository locally and will go through the source code in order to extract the different class names, method signatures and test classes of the source code. Once done it will make requests to SOCORA using the information extracted. SOCORA will response with all the components that are similar to the components extracted from the project. Simile will get this response and it will generate a report listing all the similar components to the components that are being developed by the team. This report will be sent by email to the members of team so they will know if there is already a component that can be reuse instead of developing it from scratch. Figure \ref{} describes the process of our approach.

The main objective of Simile is to help the team to find similar components so they would be able to reuse them. Thus they would reduce time development and costs through reuse components that are already tested and ready to work.

\section{Implementation}
In this section we will detail how the prototype was developed and the technologies used. First we will start explaining the Simile service which is the starting point of the whole process about cloning the project, analyse code, search components and send result. Then we will continue with the  HTTP controller which receives the Git repository of the project to be analysed, the specific branch of the repository, and the email where the result will be sent. Then we will explain the Cloner object which is in charge of cloning the project locally. After that, we will explain how the tool analyses the code and extracts the interface signatures and test classes which will be used to make the requests to SOCORA. Next we will describe how we receive the result from SOCORA. Finally we will explain how we handle the result from the previous step and build the email that will be sent to user.

\textbf{NOTE: The code snippet shown in the following subsections has been cleaned and it does not represent the source code. To see the source code with all the details, please refer to the appendix \ref{append:simile-code} or to the CD attached to this work.}

\subsection{Simile}
The Java class Simile.java [\ref{Simile.java}] is a Spring service which contains only one method, \emph{searchForComponents}. This methods receives three parameters, repo, branch, and folder. \emph{repo} represents the Git repository of the project to be analysed. \emph{branch} is the branch of the Git repository (optional). Finally \emph{folder} represents the name of the folder where the project will be cloned. Listing \ref{searchForComponents} represents the method.

\lstinputlisting[
  language=Java, numbers=left, firstnumber=1, breaklines=true, 
  basicstyle=\footnotesize,
  numberstyle=\tiny,
  caption={Method searchForComponents of Simile.java},
  captionpos=b,
  label=searchForComponents
]
{code/Simile-searchForComponents.txt}

First, we clone the project in the server (lines 3 - 5). Then we analyses the Java main code of the project cloned (lines 7 - 9) and the test classes (lines 11 - 13). When we analyse the code we extract interface signatures and test classes, we will explain how we do this in section \ref{code-analysis}. With the result of the previous step, we make the requests to SOCORA in order to get the similar components.

It is important to notice that this method is asynchronous by the Java annotation \emph{@Async}. This allows us to run this method many times in parallel.

\subsection{Cloner}
The Java class Cloner.java [\ref{Cloner.java}] is in charged of cloning the project in a local directory. The listing \ref{cloneRepository} represents the methods cloneRepository.

\lstinputlisting[
  language=Java, numbers=left, firstnumber=1, breaklines=true, 
  basicstyle=\footnotesize,
  numberstyle=\tiny,
  caption={Method cloneRepository of Cloner.java},
  captionpos=b,
  label=cloneRepository
]
{code/Cloner-cloneRepository.txt}

This method runs the command \emph{git clone} with the parameters \emph{repo}, \emph{branch}, and \emph{folder}. The command is built by the private method \emph{setCommand}.

\subsection{Interface signature and test classes extraction}
\label{code-analysis}
In this section we will explain how our prototype explores the project cloned , and how it extracts the interface signatures and the test classes from the code. 



The Java class JavaClassHandler.java [\ref{JavaClassHandler.java}] is in charge of parse the Java classes of the project to which we will search similar components for. This class is an implementation of the interface \emph{FileHandler} [\ref{FileHandler.java}]. The method implemented is \emph{handle} (listing \ref{handle}).

\lstinputlisting[
  language=Java, numbers=left, firstnumber=1, breaklines=true, 
  basicstyle=\footnotesize,
  numberstyle=\tiny,
  caption={Method handle of JavaClassHandler.java},
  captionpos=b,
  label=handle
]
{code/JavaClassHandler-handle.txt}

This method accepts three parameters: \emph{level}, \emph{path}, and \emph{file}. \emph{level} represents the level of the directory. \emph{path} represents the path of the current directory or 

\subsection{HTTP Controller}
The Java class EntryController.java [\ref{EntryController.java}] is in charge of receiving the request from a customer with the Git repository and the branch of the project to be analysed, and the email where the result of the similar component will be sent. In the listing \ref{setupRepository} we can see the method setupRepository which represents the entry point which is a POST request with the three parameters described before. The parameter \emph{repo} is required and represent the Git repository of the project that will be analysed and to which similar components will be searched for. The parameter \emph{branch} represents the branch of the git repository, this parameter is optional. Finally, the parameter \emph{email} represents the email where the result of the similar components will be sent to.

\lstinputlisting[
  language=Java, numbers=left, firstnumber=1, breaklines=true, 
  basicstyle=\footnotesize,
  numberstyle=\tiny,
  caption={Method setupRepository of EntryController.java},
  captionpos=b,
  label=setupRepository
]
{code/EntryController-setupRepository.txt}

When a request is received by the controller, it firstly validates if the parameters are correct. If they are not valid, it will respond with HTTP status code 400 and a message describing the errors. Otherwise, it triggers asynchronously the method searchForComponents with the repository, the branch, and a random CUID \footnote{CUID: Collision-resistant ids - \url{https://github.com/graphcool/cuid-java} (accessed: 04.02.2017).} as a folder name of the service Simile. This service is in charge of starting the process of cloning the project, analyse the code, make the request to SOCORA, and to send the email with the result to the customer.
\subsection{Cloner}
\subsection{Mailer}

\section{CI integration}
To integrate our prototype to a Continuous Integration process we decided to create a plugin. We decided to create our main prototype as independent as possible thus we do not depend on any CI tool. Therefore in this work to prove our approach we built a plugin for the CI server Jenkins. We decided to use Jenkins because is one of the most popular CI server for Java projects \footnote{\url{https://www.cloudbees.com/sites/default/files/2016-jenkins-community-survey-responses.pdf} (accessed: 04.02.2017)} and it is open-source.

The Simile Jenkins plugin is simple. First when a user creates a job in Jenkins, he/she can add a post build step. There he/she select the Simile plugin. When it is done, a new section is added in the task configuration (figure \ref{}). There he/she should add the Git repository of the project, the branch (optional), and the email. Once done, whenever the job is triggered, the plugin will be triggered too. The plugin sends via HTTP POST request the repository, branch, and email to Simile.

% ADD IMAGE OF JENKINS

In the future if Simile needs to be integrated to another CI server, we just need to develop a simple plugin which sends the information needed to look for similar components.
\section{SOCORA Integration}
Our prototype uses SOCORA to get similar components. After extracting the interface signatures and test classes from a project, Simile make requests to SOCORA using this information. SOCORA then searches for similar components using Merobase and ranks the candidates using non-functional requirements. Once it is done, SOCORA sends the result to Simile. 

Simile communicates with SOCORA by making HTTP requests. 
\chapter{Discussion}
\label{discussion}
The first question that arouse during this work is: \emph{Is this approach beneficial?}. As a first attempt to improve software reuse by using continuous integration, it is difficult to answer the question. In fact, it is more likely that at this version it might not be very helpful due to its immaturity. However, it is very promising in the future. In this chapter we will discuss about the possible pros and cons of this work.

This approach reduces the effort of searching suitable components by analyzing the code in background. Simile was conceived as unobtrusive tool as possible, that is why we found a great opportunity in the continuous integration process and the agile technique TDD. When a project is developed using TDD, the tests are written before the code that need to be implemented. Also according to the CI workflow, developers should commit changes as often as possible which trigger the CI server to validate the changes by testing the code. At this moment we visualized a great chance to look for similar components using textual and test-driven search approaches. Whenever the CI server is triggered we have the chance to analyze the code, extract interface signatures, and test classes which are sent to the code-search engine to retrieve similar components to the components that are being implemented. Thus developers might be able to reduce risk and development cost by reusing components which implement the functionality they are working on. Moreover it reduces the effort needed for searching components by running in background without disturbing the work of the developers.

Our prototype is able to provide high precision result that fully meet functional requirements. As test-driven search technique is able to provide the best results from the semantic point of view \cite{Reiss2009,Hummel2008,Hummel2013}, we can assure that the precision of the components returned is very high (depending on the quality of the tests). We achieve this by relying on Merobase as code-search engine.

Although Simile can provide high precision on the components retrieved, if these are not ranked accordingly the effort invested on reusing components would increase. We solve this problem by relying on SOCORA which is able to rank components using combination of non-functional properties such as performance or maintenance \cite{Kessel2016}.

The future of Simile appear to be attractive due to the massive adoption of CI/CD. Continuous Integration/Delivery adoption rates is showing an unexpected increment. According to the research made by Perforce Software \cite{Perforce2015}, 65\% of companies have adopted CI/CD in some or all projects. Moreover, 40\% of open-source projects in Github have also adopted Continuous Integration \cite{Hilton2016}. This is a great opportunity for this tool to be used vastly by companies and open-source projects, helping them to reduce development costs and time.

%%%%% CONS

This approach also has some possible cons. If a team does not use TDD or any test, Simile would only be able to extract interface signatures from the code. Therefore the result retrieved would not be as precise as it would be using tests. This is due to the fact that the search would be based on keyword and text-matching techniques which do not necessarily meet functional requirements.

Another cons is the high time it might takes to deliver a result. Test-driven search component is heavy-weight task because it requires to run tests to for each possible candidate. Although it has been improved by parallelization of test execution and optimizing the underlying code-search engine, it still takes a reasonable time for returning a result. This time can be even bigger depending on the complexity of the tests used for searching components.

One more cons is the bad quality of the documentation of the candidate components. Although the components retrieved meet the functional requirements of the user, if they are poorly documented or they are hard to understand the effort invested on reusing them will be high. Therefore it would be helpful to consider the quality of the code and documentation of the components retrieved.


 
\chapter{Future Work}
\label{future-work}
Although the approach presented in this work seems promising, there are many aspect that can be improved or need to be confirmed.

The current version need to be improved and include more features. 

It is necessary to test the approach in real projects in order to evaluate the usefulness of it. One possibility would be to introduce Simile in small open-source projects that adopt test-driven development and some CI server. Another alternative would be to adopt the tool into a company and measure the times inner-developed components are being reused by other development teams of the company.

\chapter{Conclusion}
\label{conclusion}
In this work we have presented a novel approach which is an attempt to improve the chances of reusing software by seizing the continuous integration process. As a proof-of-concept we also presented Simile, which is a tool that analyzes code, extracts interface signatures and test classes, and send them to SOCORA for searching similar components.

We started presenting the foundations of this work. First we talked about the state-of-art of software reuse, its benefits, and challenges. We also explained the different code-search engines in the market focusing on Merobase. Then we defined what ranking components means and we explained the ranking approach SOCORA. Secondly we exposed about continuous integration, its benefits, challenges, and tools. To finalize the foundation chapter we detailed microservices and its characteristics.


%\include{kapitel/grafiken}
%\include{kapitel/fazit}
%\selectlanguage{ngerman} % wenn im Inhaltsverzeichnis Appendix stehen soll, muss Engl gewaehlt sein, fuer Anhang Deutsch
\appendix

%Bibliographie
%waehle einen der folgenden 4 Eintraege
%\bibliographystyle{literatur/natdin} %DIN Style Literaturverzeichnis, comment out pagebackref
%\bibliographystyle{literatur/IEEEtran} % IEEE Style Literaturverzeichnis
%\bibliographystyle{literatur/natdinCustomized} %DIN Style Literaturverzeichnis + Punkt hinter jeder Literaturangabe -> low level config fuer Zitate: natdin.cfg im Projektordner, comment out pagebackref
%\bibliographystyle{literatur/natdinCustomizedEnglish} %DIN Style auf Englisch getrimmt mit Punkt hinter Literaturangabe -> low lovel config fuer Zitate: natdin.cfg im Projektordner, pagebackref kann damit nicht genutzt werden
%\bibliographystyle{aer} % alternativ auch apalike, aer, apalike2... s.a. http://web.reed.edu/cis/Help/LaTeX/bibtexstyles.html
\bibliographystyle{plainnat} % very nice bib style
%note on aer: does not like inbook entries
%\bibliographystyle{natdin}

\bibliography{literatur/lit} %Pfad zur bib-Datei

%appendices can be defined here, the appendix structure has to be added manually to the toc (table of contents)
%\clearpage  %toc new page
%\addcontentsline{toc}{chapter}{Appendix} %add chapter to toc
\addpart{\appendixname}
\chapter{Simile - Code}
\label{append:simile-code}
In this section we will expose the different Java classes of our prototype.

\section{de.unimannheim.informatik.swt.simile}
\lstinputlisting[
  language=Java, numbers=left, stepnumber=5, firstnumber=1, breaklines=true, 
  basicstyle=\footnotesize,
  numberstyle=\tiny,
  caption={Simile.java},
  captionpos=b,
  tabsize=2,
  showstringspaces=false,
  label=Simile.java
]
{code/main/Simile.java.txt}

\lstinputlisting[
  language=Java, numbers=left, stepnumber=5, firstnumber=1, breaklines=true, 
  basicstyle=\footnotesize,
  numberstyle=\tiny,
  caption={ServletInitializer.java},
  captionpos=b,
  tabsize=2,
  showstringspaces=false,
  label=ServletInitializer.java
]
{code/main/ServletInitializer.java.txt}

\lstinputlisting[
  language=Java, numbers=left, stepnumber=5, firstnumber=1, breaklines=true, 
  basicstyle=\footnotesize,
  numberstyle=\tiny,
  caption={SimileApplication.java},
  captionpos=b,
  tabsize=2,
  showstringspaces=false,
  label=SimileApplication.java
]
{code/main/SimileApplication.java.txt}
\subsection{de.unimannheim.informatik.swt.simile.controllers}
\lstinputlisting[
  language=Java, numbers=left, stepnumber=5, firstnumber=1, breaklines=true, 
  basicstyle=\footnotesize,
  numberstyle=\tiny,
  caption={EntryController.java},
  captionpos=b,
  tabsize=2,
  showstringspaces=false,
  label=EntryController.java
]
{code/controller/EntryController.java.txt}
\subsection{de.unimannheim.informatik.swt.simile.model}
\lstinputlisting[
  language=Java, numbers=left, stepnumber=5, firstnumber=1, breaklines=true, 
  basicstyle=\footnotesize,
  numberstyle=\tiny,
  caption={Candidate.java},
  captionpos=b,
  tabsize=2,
  showstringspaces=false,
  label=Candidate.java
]
{code/model/Candidate.java.txt}

\lstinputlisting[
  language=Java, numbers=left, stepnumber=5, firstnumber=1, breaklines=true, 
  basicstyle=\footnotesize,
  numberstyle=\tiny,
  caption={Message.java},
  captionpos=b,
  tabsize=2,
  showstringspaces=false,
  label=Message.java
]
{code/model/Message.java.txt}

\lstinputlisting[
  language=Java, numbers=left, stepnumber=5, firstnumber=1, breaklines=true, 
  basicstyle=\footnotesize,
  numberstyle=\tiny,
  caption={Metric.java},
  captionpos=b,
  tabsize=2,
  showstringspaces=false,
  label=Metric.java
]
{code/model/Metric.java.txt}
\subsection{de.unimannheim.informatik.swt.simile.services}
\lstinputlisting[
  language=Java, numbers=left, stepnumber=5, firstnumber=1, breaklines=true, 
  basicstyle=\footnotesize,
  numberstyle=\tiny,
  caption={Cloner.java},
  captionpos=b,
  tabsize=2,
  showstringspaces=false,
  label=Cloner.java
]
{code/services/Cloner.java.txt}

\lstinputlisting[
  language=Java, numbers=left, stepnumber=5, firstnumber=1, breaklines=true, 
  basicstyle=\footnotesize,
  numberstyle=\tiny,
  caption={DirectoryExplorer.java},
  captionpos=b,
  tabsize=2,
  showstringspaces=false,
  label=DirectoryExplorer.java
]
{code/services/DirectoryExplorer.java.txt}

\lstinputlisting[
  language=Java, numbers=left, stepnumber=5, firstnumber=1, breaklines=true, 
  basicstyle=\footnotesize,
  numberstyle=\tiny,
  caption={EmailSender.java},
  captionpos=b,
  tabsize=2,
  showstringspaces=false,
  label=EmailSender.java
]
{code/services/EmailSender.java.txt}

\lstinputlisting[
  language=Java, numbers=left, stepnumber=5, firstnumber=1, breaklines=true, 
  basicstyle=\footnotesize,
  numberstyle=\tiny,
  caption={JavaClassFilter.java},
  captionpos=b,
  tabsize=2,
  showstringspaces=false,
  label=JavaClassFilter.java
]
{code/services/JavaClassFilter.java.txt}

\lstinputlisting[
  language=Java, numbers=left, stepnumber=5, firstnumber=1, breaklines=true, 
  basicstyle=\footnotesize,
  numberstyle=\tiny,
  caption={JavaClassHandler.java},
  captionpos=b,
  tabsize=2,
  showstringspaces=false,
  label=JavaClassHandler.java
]
{code/services/JavaClassHandler.java.txt}

\lstinputlisting[
  language=Java, numbers=left, stepnumber=5, firstnumber=1, breaklines=true, 
  basicstyle=\footnotesize,
  numberstyle=\tiny,
  caption={JavaClassVisitor.java},
  captionpos=b,
  tabsize=2,
  showstringspaces=false,
  label=JavaClassVisitor.java
]
{code/services/JavaClassVisitor.java.txt}

\lstinputlisting[
  language=Java, numbers=left, stepnumber=5, firstnumber=1, breaklines=true, 
  basicstyle=\footnotesize,
  numberstyle=\tiny,
  caption={JavaMethodVisitor.java},
  captionpos=b,
  tabsize=2,
  showstringspaces=false,
  label=JavaMethodVisitor.java
]
{code/services/JavaMethodVisitor.java.txt}

\lstinputlisting[
  language=Java, numbers=left, stepnumber=5, firstnumber=1, breaklines=true, 
  basicstyle=\footnotesize,
  numberstyle=\tiny,
  caption={NodeIterator.java},
  captionpos=b,
  tabsize=2,
  showstringspaces=false,
  label=NodeIterator.java
]
{code/services/NodeIterator.java.txt}

\lstinputlisting[
  language=Java, numbers=left, stepnumber=5, firstnumber=1, breaklines=true, 
  basicstyle=\footnotesize,
  numberstyle=\tiny,
  caption={SocoraRequester.java},
  captionpos=b,
  tabsize=2,
  showstringspaces=false,
  label=SocoraRequester.java
]
{code/services/SocoraRequester.java.txt}
\subsection{de.unimannheim.informatik.swt.simile.util}
\lstinputlisting[
  language=Java, numbers=left, stepnumber=5, firstnumber=1, breaklines=true, 
  basicstyle=\footnotesize,
  numberstyle=\tiny,
  caption={StreamGobbler.java},
  captionpos=b,
  tabsize=2,
  showstringspaces=false,
  label=StreamGobbler.java
]
{code/util/StreamGobbler.java.txt}

\chapter{Installation of Simile}
In this chapter we will show how to install the tools we used for implementing Simile. First we will start installing the CI Server Jenkins and explain how to configure it. Then, we will continue with installing Simile Jenkins plugin. After that we will explain how to configure Simile Jenkins plugin using the test project.

\section{Installation of Jenkins}
In this section we will explain how to install Jenkins CI server using Docker and Tomcat 9.

\subsection{Docker}
For installing Jenkins, we will use Docker containers and in this subsection we will explain how to install it in different platform such as macOS and Linux.

\subsubsection{Docker on macOS}
For macOS we will use Docker for Mac tool. This is an integrated, easy-to-deploy environment for building, assembling, and shipping applications from a Mac. Moreover, it is a native Mac application architected from scratch, with a native user interface and auto-update capability, deeply integrated with OS X native virtualization, Hypervisor Framework, networking and file system, making it faster and more reliable than previous ways of getting Docker on a Mac \cite{Docker}.

First we will download the tool from the following link: \url{https://download.docker.com/mac/stable/Docker.dmg}.\\

After downloading the \textit{dmg} image, double-click on it and you will see something like figure \ref{fig:docker-mac-01}. When you get that image, drag and drop the Docker icon to the Applications folder. After that, go to Applications folder and double click on Docker application.

\begin{figure}[ht]
	\centering
    \includegraphics[width=\textwidth]{grafiken/docker-01}
    \caption{Docker for Mac Installation}
    \label{fig:docker-mac-01}
\end{figure}

\subsubsection{Docker on Linux}
In this section we will explain how to install Docker in Linux, specifically in Debian (Jessie). For other distros please refer to Docker documentation\footnote{https://docs.docker.com/engine/installation/linux/ (accessed: 21.01.2017)}.

First of all we need to update the repositories.

\begin{minted}[linenos,
               numbersep=5pt,
               frame=lines,
               framesep=2mm]{bash}
$ sudo apt-get udpate
\end{minted}

Then we need to install the packages to allow \textit{apt} to use repositories iver HTTPS.

\begin{minted}[linenos,
               numbersep=5pt,
               frame=lines,
               framesep=2mm]{bash}
$ sudo apt-get install apt-transport-https \
                       ca-certificates \
                       software-properties-common
\end{minted}

Then we add the official Docker's GPG key.

\begin{minted}[linenos,
               numbersep=5pt,
               frame=lines,
               framesep=2mm]{bash}
$ curl -fsSL https://yum.dockerproject.org/gpg | sudo apt-key add -
\end{minted}

Finally we add the stable repository of Docker and update repositories.

\begin{minted}[linenos,
               numbersep=5pt,
               frame=lines,
               framesep=2mm]{bash}
$ sudo add-apt-repository \
       "deb https://apt.dockerproject.org/repo/ \
       debian-$(lsb_release -cs) \
       main"
$ sudo apt-get update
\end{minted}

Once we added the new repositories, we update the repository list.

\begin{minted}[linenos,
               numbersep=5pt,
               frame=lines,
               framesep=2mm]{bash}
$ sudo apt-get update
\end{minted}

Then we install docker.

\begin{minted}[linenos,
               numbersep=5pt,
               frame=lines,
               framesep=2mm]{bash}
$ sudo apt-get -y install docker-engine
\end{minted}

After the installation is done we can make a little test to be sure that everything was installed correctly. The following command will download a test image and runs it in a container. Once it is running will print an informational message and exits.

\begin{minted}[linenos,
               numbersep=5pt,
               frame=lines,
               framesep=2mm]{bash}
$ sudo docker run hello-world
\end{minted}

\subsection{Jenkins in Docker}
Open a terminal on your mac and enter the following command to download the Jenkins image for Docker.

\begin{minted}[linenos,
               numbersep=5pt,
               frame=lines,
               framesep=2mm]{bash}
$ docker pull jenkins
  Using default tag: latest
  latest: Pulling from library/jenkins

  # some irrelevant output was remove
  Digest: sha256:5046434030be395ec977c98e11...
  Status: Downloaded newer image for jenkins:latest
\end{minted}

Then, we just need to run the Jenkins image with the following command.

\begin{minted}[linenos,
               numbersep=5pt,
               frame=lines,
               framesep=2mm,
               breakanywhere]{bash}
$ docker run -p 8080:8080 -p 50000:50000 jenkins
  # some irrelevant output was remove
  *************************************************************

  Jenkins initial setup is required. An admin user has been 
  created and a password generated.
  Please use the following password to proceed to installation:

  95199411f1894bfa97e937147c41aa62 # IMPORTANT, default admin pass

  This may also be found at: /var/jenkins_home/secrets/initial.
 
  *************************************************************
  # some irrelevant output was removed
  Jan 19, 2017 8:28:05 AM hudson.model.AsyncPeriodicWork$1 run
  INFO: Finished Download metadata. 25,312 ms
\end{minted}

After that, open the following URL \url{http://localhost:8080} and you should see some like figure \ref{fig:jenkins-01}.

\begin{figure}[H]
	\centering
    \includegraphics[width=0.7\textwidth]{grafiken/jenkins-01}
    \caption{First Page of Jenkins Setup}
    \label{fig:jenkins-01}
\end{figure}

There enter the password generated by Jenkins and that appeared in the logs, and then click on continue.

In the folowing page click on \textit{Install suggested plugins} like the figure \ref{fig:jenkins-02}.

\begin{figure}[H]
	\centering
    \includegraphics[width=0.8\textwidth]{grafiken/jenkins-02}
    \caption{Second Page of Jenkins Setup}
    \label{fig:jenkins-02}
\end{figure}

Then you should see something like the figure \ref{fig:jenkins-03}. Here we just need to wait untill the plugins installation finish.

\begin{figure}[H]
	\centering
    \includegraphics[width=0.8\textwidth]{grafiken/jenkins-03}
    \caption{Third Page of Jenkins Setup}
    \label{fig:jenkins-03}
\end{figure}

After the installation of plugins is done, click on \textit{Continue as admin} or create custom admin user. Then you should see the home page of Jenkins (figure )

\begin{figure}[H]
	\centering
    \includegraphics[width=0.8\textwidth]{grafiken/jenkins-04}
    \caption{Jenkins Home Page}
    \label{fig:jenkins-04}
\end{figure}

\subsection{Jenkins in Tomcat 9}
\textbf{NOTE: The following steps are based on a UNIX system.}

First, go to the following URL and download Tomcat 9: \url{http://tomcat.apache.org/download-90.cgi}.

Once downloaded, unzip the file and modify the file \emph{tomcat-users.xml} which is under the folder \emph{conf} of where Tomcat was unzipped. Add the following two lines within \emph{<tomcat-users></tomcat-users>} to setup an admin user.

\begin{minted}[linenos,
               numbersep=5pt,
               frame=lines,
               framesep=2mm,
               breakanywhere]{xml}
<role rolename="manager-gui"/>
<user username="admin" password="1234" roles="manager-gui"/>
\end{minted}

Now open a terminal and enter the following command to run Tomcat server.

\begin{minted}[linenos,
               numbersep=5pt,
               frame=lines,
               framesep=2mm,
               breakanywhere]{bash}
$ cd folder/of/tomcat
$ cd bin
$ sh catalina.sh start
\end{minted}

Then open the following URL and you should see something like Figure \ref{fig:tomcat-01}: \url{http://localhost:8080}

\begin{figure}[H]
	\centering
    \includegraphics[width=0.8\textwidth]{grafiken/tomcat-01}
    \caption{Tomcat Home Page}
    \label{fig:tomcat-01}
\end{figure}

To install Jenkins, we first need to download the \emph{.war} file. To do so go to the following URL: \url{http://mirrors.jenkins.io/war-stable/latest/jenkins.war}.

Once downloaded, move the file into the folder \emph{webapps} which is into the folder where tomcat was installed. Then enter the following commands to restart tomcat.

\begin{minted}[linenos,
               numbersep=5pt,
               frame=lines,
               framesep=2mm,
               breakanywhere]{bash}
$ sh catalina.sh stop
$ sh catalina.sh start
\end{minted}

If you go to \url{localhost:8080/jenkins} you should see the same page like Figure \ref{fig:jenkins-01}.

\section{Simile Jenkins Plugin Installation}
To install Simile Jenkins plugin, we need to follow the following steps.

First, open Jenkins URL. Then click on \textit{Manage Jenkins} (figure \ref{fig:jenkins-plugin-01}).

\begin{figure}[H]
	\centering
    \includegraphics[width=0.25\textwidth]{grafiken/jenkins-plugin-01}
    \caption{Manage Jenkins Options}
    \label{fig:jenkins-plugin-01}
\end{figure}

Then click on \textit{Manage Plugins} (figure \ref{fig:jenkins-plugin-02}).

\begin{figure}[H]
	\centering
    \includegraphics[width=0.5\textwidth]{grafiken/jenkins-plugin-02}
    \caption{Manage Plugins Options}
    \label{fig:jenkins-plugin-02}
\end{figure}

Then go to \textit{Advanced}, scroll down until \textit{Upload Plugin} section (figure \ref{fig:jenkins-plugin-03}).

\begin{figure}[H]
	\centering
    \includegraphics[width=0.9\textwidth]{grafiken/jenkins-plugin-03}
    \caption{Upload Plugin Section}
    \label{fig:jenkins-plugin-03}
\end{figure}

There click on choose file and select the hpi file \emph{simile-jenkins-plugin.hpi} provided with the CD.

\section{Simile in Tomcat 9}
To install Simile in Tomcat 9, we just need to copy the file \emph{simile.war}, provided with the CD, into the folder \emph{webapps} which into the folder where Tomcat was installed. Once done, we just need to restart tomcat.
% \newpage
% \input{kapitel/methodologyQuestionMappingTable}
% 
% \input{kapitel/question_origins}
% \include{extern/fragebogenE}
% \include{extern/fragebogenD}
% 
% \input{kapitel/invitationLetters}
% 
% 
% \include{kapitel/coverageTable}

%Eidesstattliche Erklaerung
\include{kapitel/erklaerung}

\end{document}
